\section{Results}
\label{sec:results}

\subsection{No Single Estimator Dominates the Disturbance Families}

\begin{table*}[t]
\centering
\caption{Family-level RMSE and trip-risk summary from the 34-scenario benchmark suite. RMSE values are family means of scenario means and are reported in mHz. The uncertainty term is the winner's mean within-scenario Monte Carlo standard deviation across the family.}
\label{tab:family_summary_extended}
\scriptsize
\setlength{\tabcolsep}{2.8pt}
\renewcommand{\arraystretch}{1.08}
\begin{tabular}{l c l c l c l}
\toprule
Family & Count & RMSE winner & RMSE [mHz] & Runner-up & $\Delta$ [mHz] & Lowest $T_{\mathrm{trip}}$ [s] \\
\midrule
IEEE ramps & 9 & SOGI-PLL & $19.412 \pm 6.367$ & Koopman (RK-DPMU) & 2.677 & 0 (11-way tie) \\
Magnitude steps & 7 & EKF & $0.388 \pm 0.143$ & RA-EKF & 1.451 & 0 (10-way tie) \\
Modulation & 3 & RA-EKF & $17.781 \pm 21.411$ & Koopman (RK-DPMU) & 13.659 & 0 (8-way tie) \\
Phase jumps & 3 & SOGI-FLL & $174.052 \pm 67.466$ & PLL & 4.031 & ZCD: 0.0194 \\
OOB interference & 1 & EKF & $13.641 \pm 11.132$ & RA-EKF & 7.352 & 0 (2-way tie) \\
Single-tone & 1 & ESPRIT & $0 \pm 0$ & Koopman (RK-DPMU) & 0 & 0 (13-way tie) \\
IBR harmonics & 3 & EKF & $70.464 \pm 27.297$ & TFT & 4.126 & TFT: 0.0095 \\
IBR ringdown & 5 & UKF & $313.127 \pm 23.247$ & EKF & 6.667 & ESPRIT: 0.0989 \\
IBR multi-event & 1 & PLL & $173.495 \pm 66.708$ & UKF & 2.737 & IPDFT: 0.0439 \\
Frequency step & 1 & Koopman (RK-DPMU) & $89.865 \pm 58.051$ & LKF2 & 5.291 & Koopman (RK-DPMU): 0.0099 \\
\bottomrule
\end{tabular}
\end{table*}


Table~\ref{tab:family_summary_extended} is the primary empirical object of the paper because it preserves disturbance family before any global compression takes place. The winner changes with the family. SOGI-PLL leads the nine IEEE ramp scenarios with a family-average RMSE of \SI{19.4}{mHz}. EKF leads the seven magnitude-step scenarios with \SI{0.388}{mHz} and the out-of-band interference test with \SI{13.6}{mHz}. RA-EKF leads modulation at \SI{17.8}{mHz}. SOGI-FLL gives the lowest family-average RMSE on phase jumps, UKF leads the IBR ringdown family, and PLL gives the lowest RMSE on the composite IBR multi-event case. Across the ten disturbance families, eight distinct estimators become RMSE winners. That turnover is too strong to be dismissed as noise around one latent global champion.

The clean standard families are also less informative than they first appear. IEEE ramps, magnitude steps, modulation, and the single-tone case contain many exact zero trip-risk ties. In those families, trip-risk acts mostly as a feasibility screen while RMSE carries the ranking burden. The IBR families behave differently. In IBR harmonics, IBR ringdown, and IBR multi-event stress, the lowest trip-risk method is often different from the lowest-RMSE method. The phase-jump family is the clearest example: SOGI-FLL leads RMSE, but ZCD yields the lowest mean trip-risk. Accuracy and protection risk are therefore coupled, but they are not interchangeable design objectives.

\subsection{The Hardest Cases Are Composite IBR Events and Large Phase Jumps}

The scenario-level difficulty ordering sharpens the family view. The hardest scenario in the study is the IBR multi-event case, with a median across-estimator RMSE of about \SI{825.9}{mHz}. It is followed by the IEEE and NERC \SI{60}{\degree} phase-jump cases at about \SI{675.1}{mHz} and \SI{637.5}{mHz}, then by the severe IBR ringdown-noise variants, all near or above half a hertz. These cases sit orders of magnitude above the clean single-tone and magnitude-step scenarios. The benchmark is therefore spanning distinct operating regimes rather than merely repeating one test under different random seeds.

The fragmentation of winners reinforces that point. Across the \ArchivalScenarioCount{} scenarios, ten estimators win at least one RMSE scenario outright or through exact ties. EKF is the most frequent RMSE leader with ten wins, but it is not dominant everywhere. RA-EKF and SOGI-PLL each win six scenarios, UKF and SOGI-FLL win three each, and even Koopman (RK-DPMU), ESPRIT, and Prony enter the winner set in selected cases. Scenario identity continues to matter even after family aggregation, which is exactly what a journal benchmark should surface rather than smooth away.

\begin{figure*}[!t]
  \centering
  \includegraphics[width=\textwidth,height=0.72\textheight,keepaspectratio]{archival_json_panels.pdf}
  \caption{Additional benchmark views. Top left: the hardest scenarios concentrate in composite IBR stress and large phase jumps. Top right: scenario-family difficulty spans more than three orders of magnitude. Bottom left: low RMSE and low trip-risk do not identify the same estimators. Bottom right: mean-rank consistency and compute cost define separate deployment regions.}
  \label{fig:archival_json_panels}
\end{figure*}

Figure~\ref{fig:archival_json_panels} makes those relationships more explicit. The hardest-scenario panel shows that benchmark difficulty is not spread uniformly across the catalog; it is concentrated in a small set of phase-jump and IBR stress cases. The family boxplot shows the same effect structurally: clean standard families occupy the low-error region, whereas IBR ringdown and multi-event stress define the upper tail. The lower panels show why the final recommendation cannot come from one metric alone. Estimators that look attractive on RMSE versus trip-risk are not the same estimators that dominate the mean-rank versus CPU view.

The appendix extends that compact picture into a family-resolved atlas. Figures~\ref{fig:family_dashboard_page_1} and~\ref{fig:family_dashboard_page_2} keep OOB interference, single-tone, and frequency-step baselines in separate panels rather than collapsing them into a single control bucket, and they plot the five lowest-RMSE estimators within each family across the benchmark scenarios.

\subsection{Global Summaries Help Only When Read as Secondary Evidence}

Global summaries are still useful, but only as secondary evidence. On mean rank across scenarios, TFT is the best RMSE method in the study, followed by RA-EKF and EKF. That is a different story from the family winner map because mean rank rewards repeated near-front finishes rather than dominance in the hardest families. The difference is not a contradiction; it shows that consistency and absolute error severity are different quantities.

\begin{table}[t]
\centering
\caption{Global RMSE-versus-CPU Pareto set extracted from the benchmark.}
\label{tab:pareto_set}
\scriptsize
\setlength{\tabcolsep}{3.5pt}
\renewcommand{\arraystretch}{1.08}
\begin{tabular}{l l c c}
\toprule
Estimator & Family & RMSE [mHz] & CPU [$\mu$s] \\
\midrule
RA-EKF & Model-based & 116.19 & 1.610 \\
EKF & Model-based & 123.20 & 0.813 \\
SOGI-PLL & Loop-based & 188.94 & 0.087 \\
SOGI-FLL & Loop-based & 307.87 & 0.035 \\
ZCD & Loop-based & 1.69e+05 & 0.017 \\
\bottomrule
\end{tabular}
\end{table}


The Pareto set in Table~\ref{tab:pareto_set} reveals the same distinction from a cost perspective. EKF and RA-EKF occupy the low-error middle-cost region. SOGI-PLL and SOGI-FLL remain attractive because their compute cost is at least an order of magnitude smaller. ZCD survives on the Pareto set only because its CPU cost is nearly negligible; its global error is far too large for that Pareto membership to be read as a practical recommendation. The mathematical Pareto frontier is therefore broader than the operationally credible frontier.

\begin{figure*}[!t]
  \centering
  \includegraphics[width=\textwidth,height=0.835\textheight,keepaspectratio]{Fig2_Mega_Dashboard.pdf}
  \caption{Representative views from the benchmark. The time-domain panels show why scenario family matters, while the global panels expose trade-offs between accuracy, risk, and cost that cannot be captured by a single scalar score.}
  \label{fig:dashboard_mega}
\end{figure*}

Figure~\ref{fig:dashboard_mega} is best read in that conditional way. The estimator that looks strongest in the ramp case is not the one that best handles the multi-event sequence or the noisy ringdown family, and the estimator that minimizes trip-risk is not always the one that minimizes RMSE. Once those panels are placed beside the family table and the omnibus tests, the benchmark's main claim becomes difficult to dismiss: scenario family and operating objective jointly determine what ``best'' means.
